\documentclass{exam}
\usepackage{preamble}

\pagestyle{headandfoot}
\firstpageheader{Physics}{Study Guide for Test on Unit 4: Forces}{}
\CorrectChoiceEmphasis{\color{red}\bfseries}
\SolutionEmphasis{\color{red}}

\begin{document}
\begin{questions}
\question
What is the cause of an acceleration or a change in an object's motion?

\question
In the free-body diagram shown below, which of the following is the gravitational force acting on the balloon?

\begin{figure}[h!]
    \centering
    \includegraphics[width=4cm]{Figures/Balloon.png}
\end{figure}

\question
What is the tendency of an object to resists a change in motion?

\question
The statement by Newton that for every action there is an equal but opposite reaction is which of his laws of motion?

\question
The magnitude of the force of gravity acting on an object is \fillin\ .


\question
Friction is a force that always acts \fillin\ .

\question
Compared to its mass on Earth, the mass of a 10-kg object on the moon is \fillin\ .

\question
A tennis ball and a solid steel ball with the same diameter are dropped at the same time. In the absence of air resistance, which ball has the greater acceleration?

\question
A \SI{25}{N} falling object encounters \SI{6}{N} of air resistance. The magnitude of the net force on the object is \fillin\ .

\question
How much force is needed to accelerate a 8.0-kg physics book to an acceleration of \SI{3.0}{m/s^2}?

\question
A 20-N falling object encounters \SI{20}{N} of air resistance. The magnitude of the net force on the object is \fillin\ .

\question
An object's velocity vs time graph is shown below. During which interval is there no net force acting on the object?

\begin{figure}[h!]
    \centering
    \includegraphics[width=4cm]{Figures/Velocity.png}
\end{figure}

\question
A 35.0-N physics textbook rests on a table. What is the force the table exerts on the textbook?


\clearpage
\question
Which diagram below represents a box in equilibrium?

\begin{figure}[h!]
    \centering
    \includegraphics[width=5cm]{Figures/Equilibrium.png}
\end{figure}


\question
When forces acting on an object are balanced, which quantity of motion is zero?


\question
This graph shows weight versus mass for a group of objects on planet X. What is the acceleration due to gravity on planet X?

\begin{figure}[h!]
    \centering
    \includegraphics[width=4cm]{Figures/Weight.png}
\end{figure}

\question
What is the weight of a 7.5-kilogram object on the surface of Earth?

\question
Two forces, $F_1$ and $F_2$, are applied to a block on a frictionless, horizontal surface as shown below. If the magnitude of the block's acceleration is \SI{2.0}{m/s^2}, what is the mass of the block?

\begin{figure}[h!]
    \centering
    \includegraphics{Figures/Blocks.png}
\end{figure}

\question
The diagram represents forces acting on a block that accelerates \SI{5.0}{m/s^2} on a frictionless surface. What is the mass of the block?

\begin{figure}[h!]
    \centering
    \includegraphics{Figures/Motion.png}
\end{figure}

\question
A 56-kg object is given a net force of \SI{490}{N}. What is its acceleration?


\end{questions}
\end{document}